\part{Counting}

\section{Exact State Counting}




\begin{definition}[Partitions and $q$-Pochhammer symbols]
An \emph{integer partition} of $n\geq0$ is a finite nonincreasing sequence
$\lambda=(\lambda_1,\ldots,\lambda_r)$ of positive integers with
$|\lambda|=\sum_i\lambda_i=n$.  Let $p(n)$ be the number of partitions,
with $p(0)=1$.  For a formal variable $q$, set
\[
(a;q)_n=\prod_{j=0}^{n-1}(1-aq^j),
\qquad
(a;q)_\infty=\prod_{j=0}^{\infty}(1-aq^j).
\]
Analytically, the infinite product converges for $|q|<1$.  Euler's partition
generating function is
\[
P(q)=\sum_{n\geq0}p(n)q^n=\frac1{(q;q)_\infty}.
\]
\end{definition}

\begin{exercise}[Occupation numbers]
Consider oscillator modes with one-particle energies
$\hbar\omega,2\hbar\omega,3\hbar\omega,\ldots$.  Show that a bosonic Fock
state of excitation energy $n\hbar\omega$ is a partition of $n$, and derive
\[
\prod_{m\geq1}\frac1{1-q^m}=\sum_{n\geq0}p(n)q^n.
\]
For fermionic occupation numbers in $\{0,1\}$, show that
$\prod_{m\geq1}(1+q^m)$ counts partitions into distinct parts.  Construct a
bijection, by separating powers of two in the multiplicities, between
partitions into distinct parts and partitions into odd parts, and conclude
\[
\prod_{m\geq1}(1+q^m)
=\prod_{m\geq1}\frac1{1-q^{2m-1}}.
\]
Interpret the two products as different quasiparticle descriptions with the
same exact degeneracies.
\end{exercise}

\begin{definition}[Gaussian binomial coefficients]
For $0\leq r\leq n$, the \emph{Gaussian binomial coefficient} is
\[
\qbinom{n}{r}
=\frac{(q;q)_n}{(q;q)_r(q;q)_{n-r}},
\]
and it is zero outside this range.
\end{definition}

\begin{exercise}[Gaussian binomial theorem]
Derive the recurrence
\[
\qbinom{n}{r}
=\qbinom{n-1}{r}+q^{n-r}\qbinom{n-1}{r-1}.
\]
Prove that $\qbinom{n}{r}$ is a polynomial in $q$ with nonnegative integer
coefficients and that the coefficient of $q^m$ counts partitions of $m$
whose Young diagrams fit inside an $r\times(n-r)$ rectangle.
\end{exercise}


\begin{definition}[Rogers--Ramanujan sum and product series]
For $|q|<1$, define
\[
G_{\rm sum}(q)=\sum_{r\geq0}\frac{q^{r^2}}{(q;q)_r},
\qquad
H_{\rm sum}(q)=\sum_{r\geq0}\frac{q^{r(r+1)}}{(q;q)_r},
\]
\[
G_{\rm prod}(q)=\frac1{(q;q^5)_\infty(q^4;q^5)_\infty},
\qquad
H_{\rm prod}(q)=\frac1{(q^2;q^5)_\infty(q^3;q^5)_\infty}.
\]
\end{definition}

\begin{exercise}[Rogers--Ramanujan counting]
Fix the number $r$ of parts.  Subtract the staircase
$(2r-1,2r-3,\ldots,1)$ to prove that
$q^{r^2}/(q;q)_r$ generates partitions with exactly $r$ parts whose
successive parts differ by at least two.  Use the staircase
$(2r,2r-2,\ldots,2)$ for the same difference condition with smallest part at
least two.  Interpret the product sides as partitions into parts congruent
to $1$ or $4$ modulo $5$, and to $2$ or $3$ modulo $5$, respectively.
Compute all four series through $q^{30}$ using finite products and finite
sums, and verify coefficientwise that
$G_{\rm sum}=G_{\rm prod}$ and $H_{\rm sum}=H_{\rm prod}$ to that order.
Translate the coefficient equalities into finite tests of the two partition
identities and identify which further step would be required for a proof in
all degrees.
\end{exercise}


\begin{exercise}[Partition asymptotics]
Put $q=e^{-t}$ and use the modular transformation of $\eta$ from
Exercise~\ref{ex:theta-eta-transformations} to prove
\[
P(e^{-t})
\sim\sqrt{\frac{t}{2\pi}}
\exp\!\left(\frac{\pi^2}{6t}\right)
\qquad(t\downarrow0).
\]
Apply a saddle-point argument, or an appropriate Tauberian theorem with its
hypotheses stated, to derive the Hardy--Ramanujan formula
\[
p(n)\sim\frac1{4n\sqrt3}
 \exp\!\left(\pi\sqrt{\frac{2n}{3}}\right).
\]
Interpret $\log p(n)$ as the microcanonical entropy of the bosonic modes in
the first exercise of this part.  Distinguish this asymptotic result from an
exact coefficient formula; see
\href{https://doi.org/10.1112/plms/s2-17.1.75}
{Hardy--Ramanujan, \emph{Proceedings of the London Mathematical Society}
\textbf{17} (1918)}.
\end{exercise}



\begin{definition}[The Virasoro algebra and its characters]
The \emph{Virasoro algebra} has generators $L_n$ for $n\in\Z$ and a central
element acting by $c$, with
\[
[L_m,L_n]=(m-n)L_{m+n}
 +\frac c{12}(m^3-m)\delta_{m+n,0}.
\]
A highest-weight vector $|h\rangle$ satisfies
$L_0|h\rangle=h|h\rangle$ and $L_n|h\rangle=0$ for $n>0$.
For a positive-energy module $H_a$ with finite-dimensional $L_0$
eigenspaces, its \emph{character} is
\[
\chi_a(\tau)=\Tr_{H_a}q^{L_0-c/24}.
\]
The shift by $-c/24$ is the Casimir energy produced by mapping the plane to
the cylinder.
\end{definition}

\begin{exercise}[Descendants, partitions, and null states]
Use the Virasoro commutation relations to move positive modes to the right
and apply the Poincar\'e--Birkhoff--Witt theorem to deduce the Verma-module
character
\[
\chi_{c,h}^{\rm Verma}(\tau)
=\frac{q^{h-c/24}}{(q;q)_\infty}.
\]
For the vacuum, impose $L_{-1}|0\rangle=0$ and find the resulting generic
vacuum descendant product.  Explain why additional null vectors subtract
states and turn unrestricted partition numbers into the more structured
$q$-series of minimal models.
\end{exercise}


\begin{definition}[Modular invariants]
A finite family of positive-energy sectors $H_a$ is a \emph{rational modular
character system} when its character vector carries a representation of the
modular group,
\[
\chi_a(-1/\tau)=\sum_bS_{ab}\chi_b(\tau),
\qquad
\chi_a(\tau+1)=\sum_bT_{ab}\chi_b(\tau).
\]
A full torus partition function has the form
\[
Z(\tau,\bar\tau)
=\sum_{a,b}M_{ab}\chi_a(\tau)\overline{\chi_b(\tau)}
\]
with nonnegative integer multiplicities.  It is modular invariant when
$M$ commutes with $S$ and $T$.
\end{definition}

\begin{exercise}[Ising state counting]
For the Ising theory with $c=1/2$ and primary weights
$0,1/2,1/16$, take
\[
\begin{aligned}
\chi_0&=\frac12\left(
 \sqrt{\frac{\vartheta_3}{\eta}}+
 \sqrt{\frac{\vartheta_4}{\eta}}\right),\\
\chi_{1/2}&=\frac12\left(
 \sqrt{\frac{\vartheta_3}{\eta}}-
 \sqrt{\frac{\vartheta_4}{\eta}}\right),\\
\chi_{1/16}&=\frac1{\sqrt2}
 \sqrt{\frac{\vartheta_2}{\eta}}.
\end{aligned}
\]
Expand each as $q^{h-c/24}$ times a power series and interpret its first
coefficients as exact descendant degeneracies.  Derive
\[
S=\frac12
\begin{pmatrix}
1&1&\sqrt2\\
1&1&-\sqrt2\\
\sqrt2&-\sqrt2&0
\end{pmatrix},
\qquad
T_{aa}=e^{2\pi i(h_a-c/24)}.
\]
Verify the modular invariance of
$|\chi_0|^2+|\chi_{1/2}|^2+|\chi_{1/16}|^2$. Using the Verlinde formula
\[
N_{ab}^{\ \ c}=\sum_x
\frac{S_{ax}S_{bx}\overline{S_{cx}}}{S_{0x}},
\]
recover the Ising fusion rules. Relate these towers to the critical Ising
spectroscopy exercise in Quantum Topology.
\end{exercise}

\begin{exercise}[The Cardy density of states]
Let a modular-invariant unitary CFT have discrete spectrum, unique vacuum,
and left and right central charges $c,\bar c$.  Set
$\tau=i\beta/(2\pi)$ and use $S$-invariance to show that the
high-temperature partition function is controlled by the vacuum at inverse
temperature $4\pi^2/\beta$.  A saddle-point inverse Laplace transform then
gives, for large left and right excitation numbers,
\[
\log\rho(N,\bar N)
\sim2\pi\sqrt{\frac{cN}{6}}
    +2\pi\sqrt{\frac{\bar c\bar N}{6}}.
\]
State how $c$ is replaced by
$c_{\rm eff}=c-24h_{\min}$ in a nonunitary theory.  For a nonrotating state
with $c=\bar c$ and $N=\bar N$, recover the thermal entropy used in the BTZ
exercise and hence reproduce the horizon entropy from boundary-state
counting; see
\href{https://doi.org/10.1016/0550-3213(86)90552-3}
{Cardy, \emph{Nuclear Physics B} \textbf{270} (1986)}.
\end{exercise}


\begin{definition}[Transfer matrices and integrability]
For a two-dimensional lattice model on a cylinder, a \emph{row transfer
matrix} $T_N(u)$ propagates one width-$N$ row to the next and depends on a
spectral parameter $u$.  The model is \emph{integrable} when it admits a
nontrivial commuting family
\[
[T_N(u),T_N(v)]=0
\qquad\text{for every }u,v,
\]
usually obtained from a Yang--Baxter or star--triangle relation.  Commuting
transfer matrices permit simultaneous spectral analysis, but integrability
alone does not guarantee that every finite-size degeneracy has a simple
closed formula.
\end{definition}

\begin{exercise}[Hard-hexagon counting]
In the hard-hexagon model, occupied vertices of a triangular lattice may not
be adjacent.  On an $N\times M$ torus define
\[
Z_{N,M}(z)=\sum_{I\text{ independent}}z^{|I|}.
\]
Index a row transfer matrix by cyclic binary rows and set its matrix element
to zero when two consecutive rows violate nearest-neighbor exclusion, and
otherwise to $z^{(|a|+|b|)/2}$.  Prove
$Z_{N,M}(z)=\Tr(T_N(z)^M)$ and compute the polynomial for several small
sizes.  Show that the coefficients are exact state counts, whereas the
largest eigenvalue controls the thermodynamic free energy.

Verify the algebraic identity
\[
z_c=\frac{11+5\sqrt5}{2}=\left(\frac{1+\sqrt5}{2}\right)^5,
\]
and, for the finite matrices already constructed, evaluate the largest
eigenvalue and its first two $z$-derivatives near $z_c$.  Plot their change
with $N$ and explain why finite matrices remain analytic even when their
large-size limit may develop a singularity.
\end{exercise}

\begin{exercise}[Theta-character spectra]
For a critical periodic system of circumference $L$ and propagation speed
$v$, take
\[
H_L=E_{\rm bulk}L+\frac{2\pi v}{L}
\left(L_0+\overline L_0-\frac{c+\bar c}{24}\right),
\qquad
P_L=\frac{2\pi}{L}(L_0-\overline L_0).
\]
Insert the Ising characters expressed through
$\vartheta_2,\vartheta_3,\vartheta_4$, and $\eta$ into the modular-invariant
partition function.  Prove that the coefficient of
$q^{h+N-c/24}\bar q^{\bar h+\bar N-\bar c/24}$ is the predicted
multiplicity at
\[
E-E_{\rm bulk}L
=\frac{2\pi v}{L}
\left(h+\bar h+N+\bar N-\frac{c+\bar c}{24}\right),
\quad
P=\frac{2\pi}{L}(h-\bar h+N-\bar N).
\]
Compute the first spectral gaps and degeneracies in the three Ising
sectors.  Deduce the universal gap ratios and the $L^{-1}$ Casimir-energy
correction, and state how finite-size energy levels or transfer-matrix
eigenvalues determine $c$, the primary weights, and $v$.
\end{exercise}

\section{Random Matrices}



\begin{exercise}[Hermite moments]
Let $X\sim N(0,1)$.  Complete the square and use the Gaussian integral to
prove
\[
\mathbb E[e^{tX}]=e^{t^2/2}.
\]
Define the probabilists' Hermite polynomials by
\[
e^{xt-t^2/2}=\sum_{n=0}^{\infty}\operatorname{He}_n(x)\frac{t^n}{n!}.
\]
Using two copies of the generating function, prove
\[
\mathbb E[\operatorname{He}_m(X)\operatorname{He}_n(X)]
=n!\,\delta_{mn}.
\]
Relate them to the oscillator polynomials by
$\operatorname{He}_n(x)=2^{-n/2}H_n(x/\sqrt2)$, and derive
\[
\mathbb E[X^{2n}]=\frac{(2n)!}{2^nn!},
\qquad \mathbb E[X^{2n+1}]=0.
\]
\end{exercise}


\begin{definition}[Dyson classes]
For $\lambda=(\lambda_1,\ldots,\lambda_N)$, define
\[
\Delta(\lambda)
=\det[\lambda_i^{j-1}]_{i,j=1}^N.
\]
Let $\mathbb F=\R,\C$, or $\mathbb H$ and set
\[
\beta=\dim_\R\mathbb F\in\{1,2,4\}.
\]
The Gaussian orthogonal, unitary, and symplectic ensembles, denoted GOE,
GUE, and GSE, consist of the self-adjoint matrices $H=H^*$ over
$\R,\C$, and $\mathbb H$, respectively, with density
\[
d\mathbb P_{N,\beta}(H)
=C_{N,\beta}
\exp\!\left[-\frac{\beta N}{4\sigma^2}\Tr(H^2)\right]dH.
\]
Here $H^*$ is conjugate transpose, and $dH$ is the Euclidean volume induced
by the invariant Hilbert--Schmidt inner product
$\langle X,Y\rangle_{\rm HS}=\operatorname{Re}\Tr(XY)$ on the real space
of self-adjoint matrices.
\end{definition}

\begin{exercise}[Eigenvalue--orbit Jacobian]
Use the exterior-power determinant formula to prove
\[
\Delta(\lambda)=\prod_{i<j}(\lambda_j-\lambda_i).
\]
For each Dyson class, let $G_\beta$ act on the self-adjoint matrices by
conjugation and let $T_\beta$ be the stabilizer of a regular diagonal
matrix.  Compute the Jacobian of the conjugation-orbit map
\[
(G_\beta/T_\beta)\times\R^N_{\rm reg}
\longrightarrow\operatorname{Herm}_N(\mathbb F),
\qquad
(UT_\beta,\lambda)\longmapsto
U\operatorname{diag}(\lambda)U^{-1},
\]
with respect to the invariant Hilbert--Schmidt metric, and prove
\[
dH=C\,|\Delta(\lambda)|^\beta
\prod_{i=1}^N d\lambda_i\,d\mu_{G_\beta/T_\beta},
\]
with the stabilizer and permutation factors included in $C$.  Integrate
over the conjugation orbits and derive
\[
p_{N,\beta}(\lambda)
=\frac1{Z_{N,\beta}}
\exp\!\left[-\frac{\beta N}{4\sigma^2}
\sum_{i=1}^N\lambda_i^2\right]
\prod_{i<j}|\lambda_i-\lambda_j|^\beta.
\]
\end{exercise}

\begin{definition}[Beta ensembles and Coulomb gases]
For an interval $I\subseteq\R$, a confining potential $V$, and
$\beta>0$, the \emph{beta ensemble} is
\[
d\mathbb P_{N,\beta,V}(\lambda)
=\frac1{Z_{N,\beta,V}}
\prod_{i<j}|\lambda_i-\lambda_j|^\beta
\prod_{i=1}^N e^{-\beta NV(\lambda_i)/2}\,d\lambda_i.
\]
Its \emph{Coulomb-gas Hamiltonian} is
\[
\mathcal H_N(\lambda)
=\frac N2\sum_{i=1}^NV(\lambda_i)
-\sum_{i<j}\log|\lambda_i-\lambda_j|,
\qquad
d\mathbb P_{N,\beta,V}=Z^{-1}e^{-\beta\mathcal H_N}\,d\lambda.
\]
The Jacobi beta ensemble is the Selberg law on $[0,1]^N$, and the circular
beta ensemble is
\[
\frac1{Z_{N,\beta}^{\rm circ}}
\prod_{i<j}|e^{i\theta_i}-e^{i\theta_j}|^\beta
\prod_{i=1}^N\frac{d\theta_i}{2\pi}.
\]
\end{definition}

\begin{exercise}[Jacobi beta ensemble]
For coupling $g>0$, set $\beta=2g$.  For $a,b>0$, take the $BC_N$
Calogero--Sutherland ground state in Jacobi coordinates $x_i\in(0,1)$,
\[
\Psi_0(x)=C
\prod_{i=1}^N x_i^{(a-1)/2}(1-x_i)^{(b-1)/2}
\prod_{i<j}|x_i-x_j|^{\beta/2}.
\]
Prove
\[
|\Psi_0(x)|^2
=\frac1{Z_{N,\beta}^{\rm Jac}(a,b)}
\prod_{i=1}^N x_i^{a-1}(1-x_i)^{b-1}
\prod_{i<j}|x_i-x_j|^\beta,
\]
where the Selberg integral gives
\[
Z_{N,\beta}^{\rm Jac}(a,b)
=\prod_{j=0}^{N-1}
\frac{
\Gamma(a+j\beta/2)\Gamma(b+j\beta/2)
\Gamma(1+(j+1)\beta/2)}{
\Gamma(a+b+(N+j-1)\beta/2)\Gamma(1+\beta/2)}.
\]
Interpret $|\Psi_0|^2dx_1\cdots dx_N$ as the joint law of $N$ random
eigenvalues.  For the type-$A$ circular ground state, prove instead
\[
|\Psi_0(\theta)|^2
\propto\prod_{i<j}|e^{i\theta_i}-e^{i\theta_j}|^\beta.
\]
\end{exercise}


\begin{exercise}[Semicircle equilibrium]
For the empirical measure $L_N=N^{-1}\sum_i\delta_{\lambda_i}$, derive
the mean-field energy
\[
\mathcal I_V(\mu)
=\int V(x)\,d\mu(x)
-\iint\log|x-y|\,d\mu(x)d\mu(y).
\]
Prove that its minimizer satisfies
\[
V(x)-2\int\log|x-y|\,d\mu_V(y)=\ell
\]
on its support.  For $V(x)=x^2/(2\sigma^2)$, solve the singular integral
equation and obtain
\[
\rho_{\rm sc}(x)
=\frac1{2\pi\sigma^2}
\sqrt{4\sigma^2-x^2}\,
\mathbf1_{[-2\sigma,2\sigma]}(x).
\]
\end{exercise}

\begin{definition}[Wigner matrices]
A \emph{Wigner matrix} $H_N$ is an $N\times N$ Hermitian random matrix
whose upper-triangular entries are independent and centered, with
$\mathbb E|H_{ij}|^2=\sigma^2/N$ off the diagonal and, for every fixed
$m$,
\[
\sup_{N,i,j}\mathbb E|\sqrt N H_{ij}|^m<\infty.
\]
Its \emph{empirical spectral measure} is
\[
L_{H_N}=\frac1N\sum_{j=1}^N\delta_{\lambda_j(H_N)},
\qquad
\int x^k\,dL_{H_N}(x)=\frac1N\Tr(H_N^k).
\]
\end{definition}

\begin{exercise}[Wigner semicircle law]
Expand $N^{-1}\mathbb E\Tr(H_N^k)$ as a sum over closed index walks.
Prove that only noncrossing pair-matched walks survive as $N\to\infty$
and that their number is the Catalan number.  Prove that the variance tends
to zero and conclude that $L_{H_N}$ converges in probability, in moments,
to $\rho_{\rm sc}(x)dx$.
\end{exercise}

\begin{definition}[Brownian motion]
A \emph{standard Brownian motion} is a continuous real process
\(B_t\) with \(B_0=0\) and independent increments satisfying
\[
B_t-B_s\sim N(0,t-s)\qquad(0\leq s<t).
\]
Its It\^o differentials satisfy
\[
dB_i\,dB_j=\delta_{ij}\,dt,\qquad dB_i\,dt=dt^2=0.
\]
\end{definition}

\begin{definition}[Dyson Brownian motion]
For $\beta\geq1$ and a confining potential $V$, \emph{Dyson Brownian motion}
is the interacting diffusion
\[
d\lambda_i
=\sqrt{\frac{2}{\beta N}}\,dB_i
+\left(
\frac1N\sum_{j\ne i}\frac1{\lambda_i-\lambda_j}
-\frac12V'(\lambda_i)
\right)dt.
\]
Its reversible density is
\[
Z^{-1}
\prod_{i<j}|\lambda_i-\lambda_j|^\beta
\prod_i e^{-\beta NV(\lambda_i)/2}.
\]
\end{definition}

\begin{exercise}[Eigenvalue dynamics]
Apply It\^o calculus to Hermitian and real-symmetric Ornstein--Uhlenbeck
matrix processes and derive Dyson Brownian motion for $\beta=2$ and
$\beta=1$.  Prove noncollision, compute the generator in divergence form,
and verify the stated reversible measure.  Derive the complex Burgers
equation for the large-$N$ Stieltjes transform and, in the unitary case,
the sine- and Airy-kernel local equilibria.
\end{exercise}

\begin{definition}[Unfolded spectra and level repulsion]
For a smooth approximation $\overline N(E)$ to the cumulative level count,
define
\[
x_j=\overline N(E_j),
\qquad
s_j=x_{j+1}-x_j.
\]
The variables $x_j$ are the \emph{unfolded levels}.  Their local mean
spacing is one.  Exact symmetry sectors are unfolded separately.  The
factor $|\Delta(\lambda)|^\beta$ defines \emph{level repulsion} of order
$\beta$.
\end{definition}

\begin{exercise}[Level repulsion]
For $2\times2$ GOE and GUE matrices, separate the trace from the traceless
coordinates and compute the radial distribution of the eigenvalue gap.
After rescaling to unit mean spacing, derive
\[
P_1(s)=\frac\pi2s e^{-\pi s^2/4},
\qquad
P_2(s)=\frac{32}{\pi^2}s^2e^{-4s^2/\pi}.
\]
Use the beta-ensemble density to prove
\[
P_\beta(s)\sim c_\beta s^\beta
\qquad(s\downarrow0),
\]
and compare it with the Poisson spacing law $e^{-s}$.
\end{exercise}

\begin{definition}[Global, bulk, and edge scales]
Order the eigenvalues as $\lambda_1\leq\cdots\leq\lambda_N$.  Define the
classical location $\gamma_j$ by
\[
\frac jN=\int_{-2\sigma}^{\gamma_j}\rho_{\rm sc}(x)\,dx.
\]
At $E\in(-2\sigma,2\sigma)$, the bulk coordinate is
\[
u=N\rho_{\rm sc}(E)(\lambda-E),
\]
and at the upper soft edge it is
\[
\xi=\frac{N^{2/3}}{\sigma}(\lambda-2\sigma).
\]
\end{definition}

\begin{exercise}[Bulk and edge scales]
Use the expected count $N\rho_{\rm sc}(E)dE$ to derive the bulk spacing
scale $[N\rho_{\rm sc}(E)]^{-1}$.  Prove
\[
\rho_{\rm sc}(2\sigma-s)
\sim\frac1{\pi\sigma^{3/2}}\sqrt{s}
\qquad(s\downarrow0),
\]
and derive the edge scale $s\asymp\sigma N^{-2/3}$.
\end{exercise}

\begin{definition}[Airy kernel and the Tracy--Widom law]
Let $\operatorname{Ai}$ be the solution of $y''=xy$ decaying at
$+\infty$.  Define
\[
K_{\rm Ai}(x,y)
=\frac{\operatorname{Ai}(x)\operatorname{Ai}'(y)
-\operatorname{Ai}'(x)\operatorname{Ai}(y)}{x-y}
\]
by continuity on the diagonal.  The GUE Tracy--Widom distribution is
\[
F_2(s)=\det(I-K_{\rm Ai})_{L^2(s,\infty)}.
\]
\end{definition}


\begin{definition}[Point processes and correlation functions]
Let $(X,\mu)$ be a locally compact measure space.  A \emph{simple point
process} is a random locally finite counting measure
\[
\Xi=\sum_i\delta_{x_i}
\]
with distinct points almost surely.  Its $k$-point correlation function
$\rho_k$ is defined by
\[
\mathbb E\sum_{i_1,\ldots,i_k}^{\ne}
f(x_{i_1},\ldots,x_{i_k})
=\int_{X^k}f(x_1,\ldots,x_k)\rho_k(x_1,\ldots,x_k)
\prod_{j=1}^k d\mu(x_j).
\]
For a symmetric joint density $p_N$ of exactly $N$ points,
\[
\rho_k(x_1,\ldots,x_k)
=\frac{N!}{(N-k)!}
\int_{X^{N-k}}p_N(x_1,\ldots,x_N)
\prod_{j=k+1}^N d\mu(x_j).
\]
\end{definition}

\begin{exercise}[Factorial moments and gap probabilities]
For a relatively compact measurable set $B$ and $N_B=\Xi(B)$, prove,
whenever either side is finite,
\[
\mathbb E(N_B)_k=\int_{B^k}\rho_k\,d\mu^{\otimes k},
\qquad
(N_B)_k=N_B(N_B-1)\cdots(N_B-k+1).
\]
If $\mathbb E[2^{N_B}]<\infty$, use inclusion--exclusion and justify the
exchange of expectation and summation to prove
\[
\mathbb P(N_B=0)
=\sum_{k=0}^{\infty}\frac{(-1)^k}{k!}
\int_{B^k}\rho_k\,d\mu^{\otimes k}.
\]
\end{exercise}


\begin{definition}[Determinantal and Pfaffian processes]
For an antisymmetric matrix $A=(a_{ij})\in M_{2m}(\C)$, define
\[
\operatorname{Pf}(A)
=\frac1{2^m m!}
\sum_{\sigma\in S_{2m}}\operatorname{sgn}(\sigma)
\prod_{j=1}^m a_{\sigma(2j-1),\sigma(2j)}.
\]
A point process is \emph{determinantal} with kernel $K$ when
\[
\rho_k(x_1,\ldots,x_k)
=\det[K(x_i,x_j)]_{i,j=1}^k.
\]
It is \emph{Pfaffian} with antisymmetric $2\times2$ matrix kernel
$\mathbb K$ when
\[
\rho_k(x_1,\ldots,x_k)
=\operatorname{Pf}[\mathbb K(x_i,x_j)]_{i,j=1}^k.
\]
The Fredholm Pfaffian is the corresponding Pfaffian analogue of the
Fredholm determinant.  A \emph{positive contraction} is an
operator $K$ satisfying $0\leq K\leq I$.
\end{definition}

\begin{exercise}[Counting statistics and Fredholm formulas]
For a determinantal process, prove
\[
\mathbb E[z^{N_B}]
=\det(I+(z-1)K_B),
\qquad
\mathbb P(N_B=0)=\det(I-K_B).
\]
Derive the corresponding Fredholm-Pfaffian formula for a Pfaffian
process.  If $K_B$ is a positive trace-class contraction with eigenvalues
$(\lambda_j)$, prove
\[
\det(I+zK_B)=\prod_j(1+z\lambda_j)
\]
and that $N_B$ is distributed as a sum of independent Bernoulli variables
with parameters $\lambda_j$.
\end{exercise}


\begin{definition}[Orthogonal-polynomial ensembles]
Let $w(x)d\mu(x)$ be a measure with finite moments and let $p_j$ be its
monic orthogonal polynomials,
\[
\int p_j(x)p_k(x)w(x)d\mu(x)=h_j\delta_{jk}.
\]
The associated \emph{orthogonal-polynomial ensemble} has density
\[
p_N(x_1,\ldots,x_N)
=\frac1{Z_N}\Delta(x)^2\prod_{i=1}^Nw(x_i),
\qquad
Z_N=N!\prod_{j=0}^{N-1}h_j.
\]
Its correlation kernel is
\[
K_N(x,y)
=\sqrt{w(x)w(y)}
\sum_{j=0}^{N-1}\frac{p_j(x)p_j(y)}{h_j}.
\]
\end{definition}

\begin{exercise}[Andr\'eief identity]
Prove Andr\'eief's identity
\[
\int_{X^N}\det[f_i(x_j)]\det[g_i(x_j)]
\prod_{j=1}^N d\mu(x_j)
=N!\det\!\left[\int_X f_i(x)g_j(x)d\mu(x)\right].
\]
Use the Vandermonde determinant to prove
\[
p_N(x_1,\ldots,x_N)
=\frac1{N!}\det[K_N(x_i,x_j)]_{i,j=1}^N
\]
and
\[
\rho_k(x_1,\ldots,x_k)
=\det[K_N(x_i,x_j)]_{i,j=1}^k.
\]
\end{exercise}

\begin{exercise}[Skew-orthogonal ensembles]
For an antisymmetric bilinear form $\langle\cdot,\cdot\rangle_s$, construct
skew-orthogonal polynomials satisfying
\[
\langle p_{2j},p_{2k+1}\rangle_s=r_j\delta_{jk},
\qquad
\langle p_{2j},p_{2k}\rangle_s
=\langle p_{2j+1},p_{2k+1}\rangle_s=0.
\]
Prove de Bruijn's Pfaffian integration identity and use it to derive the
matrix correlation kernels for the GOE and GSE.
\end{exercise}


\begin{definition}[Christoffel--Darboux kernels]
Let $q_j=p_j/\sqrt{h_j}$ and let $\kappa_j$ be the leading coefficient of
$q_j$.  The Christoffel--Darboux kernel is
\[
K_N(x,y)
=\sqrt{w(x)w(y)}\,\frac{\kappa_{N-1}}{\kappa_N}
\frac{q_N(x)q_{N-1}(y)-q_{N-1}(x)q_N(y)}{x-y}.
\]
An \emph{integrable kernel} is a kernel of the form
\[
K(x,y)=\frac{f(x)^{\mathsf T}g(y)}{x-y},
\qquad
f(x)^{\mathsf T}g(x)=0.
\]
\end{definition}

\begin{exercise}[Projection and resolvent identities]
Derive the Christoffel--Darboux formula from the three-term recurrence.
Prove
\[
\int K_N(x,z)K_N(z,y)d\mu(z)=K_N(x,y),
\qquad
\int K_N(x,x)d\mu(x)=N.
\]
For an integrable trace-class kernel $K$, prove that the resolvent
$R=K(I-K)^{-1}$ is again integrable and express its numerator using
$(I-K)^{-1}f$ and $(I-K^{\mathsf T})^{-1}g$.
\end{exercise}


\begin{definition}[Universal limiting kernels]
The bulk, soft-edge, and hard-edge kernels are respectively
\[
K_{\sin}(u,v)=\frac{\sin\pi(u-v)}{\pi(u-v)},
\]
\[
K_{\rm Ai}(u,v)
=\frac{\operatorname{Ai}(u)\operatorname{Ai}'(v)
-\operatorname{Ai}'(u)\operatorname{Ai}(v)}{u-v},
\]
and, for $\alpha>-1$,
\[
K_{\rm Bes}^{(\alpha)}(u,v)
=\frac{
J_\alpha(\sqrt u)\sqrt v\,J_\alpha'(\sqrt v)
-\sqrt u\,J_\alpha'(\sqrt u)J_\alpha(\sqrt v)}
{2(u-v)}.
\]
The values on the diagonal are defined by continuity.
\end{definition}

\begin{exercise}[Classical kernel universality]
Use the Christoffel--Darboux formula and the asymptotics of the Hermite,
Laguerre, and Jacobi polynomials to prove convergence of their rescaled
kernels, where applicable, to $K_{\sin}$ in the bulk, $K_{\rm Ai}$ at a
soft edge, and $K_{\rm Bes}^{(\alpha)}$ at a hard edge.  Deduce convergence
of every correlation function and of the associated Fredholm gap
probabilities.  Recover the Tracy--Widom law for the largest GUE eigenvalue.
For an unfolded spectrum, derive from $K_{\sin}$ the nearest-neighbor
level-repulsion exponent and number variance, and state the corresponding
predictions for a chaotic quantum spectrum in the unitary symmetry class.
From $K_{\rm Ai}$ derive the
$N^{-2/3}$ edge-fluctuation scale and the limiting distribution of the
largest resonance or covariance eigenvalue.  Specify the unfolding and
centering required before comparing either prediction with spectral data.
\end{exercise}

\section{Free Probability}



\begin{definition}[Noncommutative probability spaces]
A \emph{noncommutative probability space} is a unital complex algebra
$A$ with a unital linear functional $\varphi:A\to\C$.  Its elements are
\emph{random variables} and their joint moments are
$\varphi(a_1\cdots a_n)$.  In the $*$-algebra setting, a \emph{state} is a
positive unital functional; it is \emph{tracial} when
$\varphi(ab)=\varphi(ba)$.
\end{definition}

\begin{definition}[Free independence and free cumulants]
Unital subalgebras $(A_i)_i$ of $(A,\varphi)$ are \emph{freely independent}
if
\[
\varphi(a_1\cdots a_m)=0
\]
whenever $a_j\in A_{i_j}$, $\varphi(a_j)=0$, and neighboring indices are
distinct. Variables are free when their generated unital subalgebras are.

Let $NC(n)$ be the lattice of noncrossing partitions of
$\{1,\ldots,n\}$. The \emph{free cumulants} $\kappa_n$ are the multilinear
functionals determined recursively by
\[
\varphi(a_1\cdots a_n)
=\sum_{\pi\in NC(n)}\kappa_\pi(a_1,\ldots,a_n),
\]
where $\kappa_\pi$ is the product of the cumulants associated to the blocks
of $\pi$. Freeness is equivalently the vanishing of every mixed free
cumulant.
\end{definition}

\begin{definition}[Semicircular variables]
A self-adjoint variable $S$ is \emph{semicircular} with variance $\sigma^2$
if its distribution is
\[
d\mu_{\mathrm{sc},\sigma}(x)
=\frac{1}{2\pi\sigma^2}
 \sqrt{4\sigma^2-x^2}\,
 \mathbf 1_{[-2\sigma,2\sigma]}(x)\,dx.
\]
Equivalently, its only nonzero free cumulant is
$\kappa_2(S,S)=\sigma^2$.
\end{definition}

\begin{exercise}[Free central limit theorem]
Use noncrossing pairings to prove
\[
\varphi(S^{2m+1})=0,
\qquad
\varphi(S^{2m})=C_m\sigma^{2m},
\qquad
C_m=\frac1{m+1}\binom{2m}{m}.
\]
Let $X_1,X_2,\ldots$ be freely independent, identically distributed,
centered self-adjoint variables with $\varphi(X_i^2)=\sigma^2$ and moments of
every order. Use
multilinearity and vanishing mixed cumulants to prove that the moments of
$n^{-1/2}(X_1+\cdots+X_n)$ converge to these semicircular moments. Compare
this free central limit theorem with the Gaussian limit for classically
independent variables.
\end{exercise}

\begin{definition}[Cauchy and $R$-transforms]
For a compactly supported probability measure $\mu$ on $\R$, its
\emph{Cauchy transform} is
\[
G_\mu(z)=\int_\R\frac{d\mu(t)}{z-t},
\qquad z\in\C\setminus\supp\mu.
\]
Near infinity it has a local compositional inverse
$K_\mu(w)=G_\mu^{\langle-1\rangle}(w)$ near $w=0$. The
\emph{$R$-transform} is defined by
\[
K_\mu(w)=\frac1w+R_\mu(w),
\qquad
R_\mu(w)=\sum_{n\geq1}\kappa_n(\mu)w^{n-1}.
\]
If freely independent self-adjoint variables $X,Y$ have laws $\mu,\nu$,
then the law of $X+Y$ is their \emph{free additive convolution}
$\mu\boxplus\nu$, characterized by
\[
R_{\mu\boxplus\nu}=R_\mu+R_\nu.
\]
\end{definition}

\begin{exercise}[Free convolution]
Derive the additivity of the $R$-transform from the vanishing of mixed free
cumulants. Show that a centered semicircular law of variance $\sigma^2$ has
$R(w)=\sigma^2w$ and hence
\[
\sigma^2G(z)^2-zG(z)+1=0.
\]
Select the solution satisfying $zG(z)\to1$ at infinity and recover the
semicircle density from the Stieltjes inversion formula. Deduce that the free
sum of semicircular variables of variances $\sigma_1^2$ and $\sigma_2^2$ is
semicircular of variance $\sigma_1^2+\sigma_2^2$, and rederive the free
central limit theorem at the level of transforms.
\end{exercise}


\begin{definition}[Asymptotic freeness]
Random matrix families
$A_N^{(1)},\ldots,A_N^{(r)}$ are \emph{asymptotically free} when their
normalized expected joint trace moments converge to the moments of freely
independent variables:
\[
\lim_{N\to\infty}
\frac1N\mathbb E\Tr
P(A_N^{(1)},\ldots,A_N^{(r)})
=\varphi(P(a_1,\ldots,a_r))
\]
for every noncommutative polynomial $P$.
\end{definition}

\begin{exercise}[GUE asymptotic freeness]
For independent GUE matrices $H_N^{(1)},\ldots,H_N^{(r)}$, apply Wick's
rule to normalized mixed traces of centered polynomials.  Prove that
crossing pairings are lower order and that mixed free cumulants vanish in
the limit.  Conclude that the family converges in joint moments to freely
independent semicircular variables.
\end{exercise}

\begin{exercise}[Catalan semicircle transform]
Let $C_n=(n+1)^{-1}\binom{2n}{n}$.  Decompose a noncrossing pairing at the
partner of its first point to prove
\[
C_{n+1}=\sum_{j=0}^nC_jC_{n-j}.
\]
Deduce
\[
\sum_{n=0}^{\infty}C_nt^n
=\frac{1-\sqrt{1-4t}}{2t}.
\]
Use the semicircular moment formula to recover
\[
G_{\rm sc}(z)=\frac{z-\sqrt{z^2-4}}2,
\qquad zG_{\rm sc}(z)\longrightarrow1
\quad(z\to\infty),
\]
and identify the branch cut with the limiting spectral interval of GUE.
Use Stieltjes inversion to recover the semicircle density from the jump of
$G_{\rm sc}$ across the branch cut.  For a measured large Hamiltonian or
coupled-mode spectrum modeled by an additive free sum, use
$R_{\mu\boxplus\nu}=R_\mu+R_\nu$ to predict the support edges and density
of states.  State how convergence of the empirical Stieltjes transform and
the observed spectral edges test the asymptotic-freeness model.
\end{exercise}

\section{Integrable Probability}



\begin{definition}[Tableaux and Schur functions]
A semistandard Young tableau of shape $\lambda$ has weakly increasing rows
and strictly increasing columns.
For variables $a=(a_1,\ldots,a_m)$, define
\[
s_\lambda(a)=\sum_T\prod_{i=1}^m a_i^{m_i(T)},
\]
where $m_i(T)$ is the number of entries equal to $i$.
\end{definition}

\begin{definition}[RSK correspondence]
The Robinson--Schensted--Knuth construction associates
\[
A=(a_{ij})\longmapsto(P,Q)
\]
to a matrix of nonnegative integers a pair of semistandard Young
tableaux of the same shape $\lambda$, with the row and column sums of
$A$ giving the contents of $P$ and $Q$.  For weights $a_{ij}\geq0$, define
\[
G(m,n)=\max_{\pi:(1,1)\nearrow(m,n)}
\sum_{(i,j)\in\pi}a_{ij}.
\]
\end{definition}

\begin{exercise}[Schensted and Greene theorems]
Prove the RSK bijection and Greene's theorem.  For
$A\in M_{m,n}(\Z_{\geq0})$, deduce
\[
\lambda_1=G(m,n)
\]
and identify $\lambda_1+\cdots+\lambda_r$ with the maximal weight of
$r$ disjoint up-right paths.  For independent geometric weights with
parameters $a_i b_j$, use RSK to derive
\[
\mathbb P(\lambda)
=\prod_{i,j}(1-a_i b_j)s_\lambda(a)s_\lambda(b)
\]
and prove the Cauchy identity
\[
\sum_\lambda s_\lambda(a)s_\lambda(b)
=\prod_{i,j}(1-a_i b_j)^{-1}.
\]
\end{exercise}


\begin{definition}[Schur measures and processes]
For $\mu\subseteq\lambda$, the \emph{skew Schur function}
$s_{\lambda/\mu}$ is the tableau sum over the skew diagram
$\lambda\setminus\mu$.  A \emph{specialization} is an algebra
homomorphism from the algebra of symmetric functions to $\C$; it is
\emph{Schur-positive} when
$s_{\lambda/\mu}(\rho)\geq0$ for all $\mu\subseteq\lambda$.
The \emph{Schur measure} associated to two Schur-positive specializations
$\rho^+,\rho^-$ is
\[
\mathbb P(\lambda)
=\frac{s_\lambda(\rho^+)s_\lambda(\rho^-)}
{\sum_\nu s_\nu(\rho^+)s_\nu(\rho^-)}.
\]
An ascending \emph{Schur process} is a probability measure on
$\lambda^{(1)}\subseteq\cdots\subseteq\lambda^{(m)}$ proportional to
\[
s_{\lambda^{(1)}}(\rho_1^+)
\prod_{r=2}^m
s_{\lambda^{(r)}/\lambda^{(r-1)}}(\rho_r^+)
s_{\lambda^{(m)}}(\rho^-).
\]
\end{definition}

\begin{exercise}[Schur-process correlations]
Encode a partition by the Maya diagram
\[
\mathcal X(\lambda)=\{\lambda_i-i+\tfrac12:i\geq1\}.
\]
Use the Cauchy and skew-Cauchy identities to prove that the induced
space--time point process is determinantal.  Derive its double-contour
correlation kernel and express the distribution of $\lambda_1$ as a
Fredholm determinant.  Recover geometric last-passage percolation as a
one-time marginal.
\end{exercise}

\begin{definition}[Macdonald measures and processes]
Let $P_\lambda(\,\cdot\,;q,t)$ be the Macdonald polynomials and let
$Q_\lambda=b_\lambda(q,t)P_\lambda$ be their dual family.  Define the skew
functions by
\[
P_\lambda(x,y)=\sum_\mu P_{\lambda/\mu}(x)P_\mu(y),
\qquad
Q_\lambda(x,y)=\sum_\mu Q_{\lambda/\mu}(x)Q_\mu(y).
\]
For finite
specializations $a,b$, define
\[
\Pi(a;b)
=\prod_{i,j}\frac{(t a_i b_j;q)_\infty}{(a_i b_j;q)_\infty}.
\]
The \emph{Macdonald measure} is
\[
\mathbb P_{q,t}(\lambda)
=\frac{P_\lambda(a;q,t)Q_\lambda(b;q,t)}{\Pi(a;b)}.
\]
The ascending \emph{Macdonald process} is
\[
\frac{
P_{\lambda^{(1)}}(a^{(1)})
\prod_{r=2}^m
P_{\lambda^{(r)}/\lambda^{(r-1)}}(a^{(r)})
Q_{\lambda^{(m)}}(b)}
{\Pi(a^{(1)},\ldots,a^{(m)};b)}.
\]
\end{definition}

\begin{exercise}[Macdonald observables]
Apply the commuting Macdonald difference operators to the Cauchy identity
and derive contour-integral formulas for the joint moments of
\[
\sum_i q^{\lambda_i}t^{N-i}.
\]
Prove that $q=t$ gives the Schur process, $q=0$ gives the
Hall--Littlewood process, and $t=0$ gives the $q$-Whittaker process.
Obtain Jack processes from the limit $q=t^\alpha$, $t\uparrow1$.
\end{exercise}


\begin{definition}[White noise]
\emph{Space--time white noise} on $\R_+\times\R$ is the centered Gaussian
generalized process $\xi$ satisfying
\[
\mathbb E[\xi(\varphi)\xi(\psi)]
=\int_{\R_+\times\R}\varphi(t,x)\psi(t,x)\,dt\,dx
\]
for test functions $\varphi,\psi$.
\end{definition}

\begin{definition}[TASEP and the KPZ equation]
In the totally asymmetric simple exclusion process, particles on $\Z$
attempt rate-one jumps from $x$ to $x+1$, and a jump is suppressed when
$x+1$ is occupied.  Its integrated current defines a height function.
The one-dimensional KPZ equation is
\[
\partial_t h
=\nu\,\partial_x^2h
+\frac{\lambda}{2}(\partial_xh)^2
+\sqrt D\,\xi,
\]
where $\xi$ is space--time white noise.  Its Cole--Hopf solution is
$h=(2\nu/\lambda)\log Z$, where $Z$ solves the corresponding stochastic
heat equation with multiplicative noise.
\end{definition}

\begin{exercise}[KPZ scaling]
Use the recursion
\[
G(m,n)=\max\{G(m-1,n),G(m,n-1)\}+a_{mn}
\]
to identify exponential last-passage times with TASEP currents.  For
i.i.d.\ exponential weights, prove
\[
\frac{G(n,n)-4n}{2^{4/3}n^{1/3}}
\Longrightarrow\chi_2,
\qquad
\mathbb P(\chi_2\leq s)=F_2(s).
\]
Derive the $1{:}2{:}3$ KPZ fluctuation scaling and its spatial Airy
process limit.
\end{exercise}

\begin{exercise}[Macdonald--KPZ degeneration]
At $t=0$, intertwine the $q$-Whittaker difference operators with the
$q$-TASEP generator.  Derive a Fredholm-determinant formula for the
$q$-Laplace transform of a particle position.  Take the $q\uparrow1$ limit
to the O'Connell--Yor semi-discrete polymer and then its
intermediate-disorder limit to the stochastic heat equation and the KPZ
equation.  Identify the corresponding degenerations of the
Macdonald-process observables.
Under narrow-wedge initial data, carry the steepest-descent analysis of
the limiting Fredholm kernel to the Airy kernel.  After fixing the
nonuniversal velocity and scale from the first two cumulants, prove that
the centered height has $t^{1/3}$ fluctuations and converges to the GUE
Tracy--Widom distribution.  Deduce the $t^{2/3}$ correlation-length scale
and state the resulting data collapse for repeated interface-height or
directed-polymer free-energy measurements.
\end{exercise}
